\documentclass[letterpaper,10pt]{article}

\usepackage[T1]{fontenc}
\usepackage{lmodern}
\usepackage{microtype}
\usepackage[letterpaper,top=0.43in,bottom=0.43in,left=0.52in,right=0.52in]{geometry}
\usepackage{array}
\usepackage{enumitem}
\usepackage{tabularx}
\usepackage{titlesec}
\usepackage{xcolor}
\usepackage[hidelinks]{hyperref}

\input{glyphtounicode}
\pdfgentounicode=1

\definecolor{accent}{HTML}{1E5BA8}
\definecolor{muted}{HTML}{4A4A5A}

\pagestyle{empty}
\setlength{\parindent}{0pt}
\setlength{\tabcolsep}{0pt}
\setlength{\emergencystretch}{2em}
\urlstyle{same}
\hypersetup{
  pdftitle={Linji Wang - AI Robotics and Systems Engineer - Curriculum Learning and Embodied RL},
  pdfauthor={Linji Wang},
  pdfsubject={Adaptive robot learning, policy integration, and systems engineering},
  pdfkeywords={AI robotics, robot learning, curriculum learning, reinforcement learning, embodied AI, Aurora PostgreSQL, database systems, query optimization, AWS}
}

\titleformat{\section}
  {\large\bfseries\color{accent}}
  {}{0pt}{}[\vspace{-4pt}\color{accent}\titlerule]
\titlespacing*{\section}{0pt}{6pt}{3pt}

\setlist[itemize]{leftmargin=0.17in,itemsep=1.2pt,topsep=1.5pt,parsep=0pt,partopsep=0pt}

\newcommand{\role}[4]{%
  \textbf{#1}\hfill\textbf{#2}\\[-1pt]
  \textit{#3}\hfill\textit{#4}\par
}

\newcommand{\education}[4]{%
  \textbf{#1}\hfill #2\\[-1pt]
  \textit{#3}\hfill\textit{#4}\par
}

\begin{document}

\begin{center}
  {\fontsize{19}{21}\selectfont\bfseries LINJI (JOEY) WANG}\\[-1pt]
  {\large\bfseries\color{accent}Robot Learning Researcher \textbar{} Robotics Systems Engineer}\\[2pt]
  \small Fairfax, VA \textbar{}
  \href{mailto:joewwang@outlook.com}{joewwang@outlook.com} \textbar{}
  412-888-6071 \textbar{}
  \href{https://linjiw.github.io}{linjiw.github.io} \textbar{}
  \href{https://www.linkedin.com/in/linjiw}{LinkedIn} \textbar{}
  \href{https://github.com/linjiw}{GitHub} \textbar{}
  \href{https://scholar.google.com/citations?user=VURUgFMAAAAJ}{Google Scholar}
\end{center}
\vspace{-7pt}

\section{PROFILE}
\small Computer Science Ph.D. researcher connecting adaptive RL training with robot policy integration. Research spans curricula over tasks, auxiliary rewards, and domain randomization, with first/co-first-author papers at IROS 2025. Engineering experience includes C++/ROS 2/ONNX inference and C/Python systems at AWS.

\section{TECHNICAL EXPERTISE}
\small
\begin{tabularx}{\textwidth}{>{\bfseries}p{1.36in} X}
Robot learning & Automatic curriculum learning, deep reinforcement learning, PPO, teacher--student learning, reward adaptation, latent representations, domain randomization, sim-to-real \\
Robotics systems & PyTorch, Isaac Gym, Isaac Lab, MuJoCo, ROS 2, ONNX Runtime; navigation, quadruped locomotion, off-road mobility, humanoid policy inference \\
Systems \& software & C, C++, Python, Aurora PostgreSQL, query processing and join optimization, AWS, Docker, Git/CI-CD \\
\end{tabularx}

\section{RESEARCH \& ENGINEERING EXPERIENCE}
\small
\role{RobotiXX Lab, George Mason University}{Fairfax, VA}{Graduate Research Assistant --- Robot Learning}{Aug 2023 -- Present}
\begin{itemize}
  \item \textbf{Task curricula --- GACL:} Developed VAE task representations, performance-history conditioning, and grounded task sampling for PPO in 128 parallel Isaac Gym environments. Reported navigation / quadruped simulation success: 81.85\% / 79.21\% versus CLUTR at 76.67\% / 74.65\% (first author, IROS 2025).
  \item \textbf{Reward curricula --- Reward Training Wheels:} Co-developed adaptive auxiliary-reward weighting. Reported off-road simulation success: 76.67\% versus 34.44\% with expert rewards; physical completion: 5/5 versus 2/5 trials (co-first author, IROS 2025).
  \item \textbf{Domain-randomization curricula --- \href{https://linjiw.github.io/research/\#lucid}{LUCID} (preprint):} Execution-informed humanoid training with a frozen temporal encoder and bounded feedback scheduler. Study reports 88.9\% versus 76.8\% full-randomization simulation completion against filtered-error PI; G1 at 40 ms added delay: 38/60 versus 23/60 trials.
  \item \textbf{Policy integration:} Implemented residual/base ONNX inference, metadata-driven observations, normalization, and temporal history in an upstream C++/ROS 2 controller for Unitree G1; developing Isaac Lab--MuJoCo parity diagnostics. This integration remains a simulation prototype.
  \item Co-authored \textbf{Moving Through Clutter} (3rd author, preprint 2026), \textbf{DDP} (3rd, IROS 2025), and \textbf{Adaptive Dynamics Planning} (4th, ICRA 2026). The DDP-based RobotiXX team placed 2nd in both 2025 BARN Challenge phases.
\end{itemize}

\role{Amazon Web Services (AWS)}{}{Software Development Engineer Intern --- Amazon RDS \& Aurora}{Summers 2025 \& 2026}
\begin{itemize}
  \item \textbf{Aurora PostgreSQL (2026):} Implemented and extended Adaptive Join in C, adding support for additional join types; built join benchmarks and synthetic-data generators to investigate regressions across engine changes.
  \item \textbf{RDS Proxy (2025):} Built Python/Streamlit performance tooling with multi-region CloudWatch metrics; implemented Welch's tests, power analysis, and Bonferroni correction for regression investigation.
\end{itemize}

\role{Computational Engineering and Robotics Lab, Carnegie Mellon University}{Pittsburgh, PA}{Research Assistant --- 3D Perception and AR}{Jan 2022 -- May 2023}
\begin{itemize}
  \item Built an AR scene-inpainting pipeline using GAN image completion plus RANSAC/DBSCAN point-cloud segmentation.
\end{itemize}

\section{SELECTED PUBLICATIONS}
\small
\begin{itemize}
  \item \textbf{Wang, L.}, Xu, Z., Stone, P., Xiao, X. ``\href{https://doi.org/10.1109/IROS60139.2025.11246910}{GACL: Grounded Adaptive Curriculum Learning with Active Task and Performance Monitoring}.'' \textit{IROS 2025}. First author.
  \item \textbf{Wang, L.*}, Xu, T.*, Lu, Y., Xiao, X. ``\href{https://doi.org/10.1109/IROS60139.2025.11247039}{Reward Training Wheels: Adaptive Auxiliary Rewards for Robotics Reinforcement Learning}.'' \textit{IROS 2025}. Co-first author.
  \item \textbf{Additional robotics publications:} Moving Through Clutter (3rd author, arXiv 2026); DDP (3rd, IROS 2025); Adaptive Dynamics Planning (4th, ICRA 2026); ColorMap-VIO / II-NVM (5th, RA-L 2026 / 2025).
\end{itemize}

\section{EDUCATION}
\small
\education{George Mason University}{Fairfax, VA}{Ph.D. in Computer Science --- Advisor: Xuesu Xiao}{2023 -- Expected 2027}
\education{Carnegie Mellon University}{Pittsburgh, PA}{M.S. in Mechanical Engineering, GPA: 3.94/4.0}{2021 -- 2023}
\education{University of Cincinnati}{Cincinnati, OH}{B.S. in Mechanical Engineering, magna cum laude}{2016 -- 2021}

\end{document}
